\documentclass[12pt,a4paper]{article}
\usepackage[utf8]{inputenc}
\usepackage{amsmath,amssymb,amsthm}
\usepackage{algorithm}
\usepackage{algorithmic}
\usepackage{geometry}
\geometry{margin=1in}

\newtheorem{theorem}{Theorem}
\newtheorem{lemma}[theorem]{Lemma}
\newtheorem{proposition}[theorem]{Proposition}
\theoremstyle{definition}
\newtheorem{definition}[theorem]{Definition}
\theoremstyle{remark}
\newtheorem{remark}[theorem]{Remark}

\title{Problem 27: Quadratic Primes}
\author{}
\date{}

\begin{document}
\maketitle

\section{Problem Statement}

Find the product $a \cdot b$ for the quadratic $n^2 + an + b$, where $|a| < 1000$ and $|b| \leq 1000$, that produces the maximum number of primes for consecutive values of~$n$ starting with $n = 0$.

\section{Mathematical Development}

\begin{definition}
For integers $a, b$, define $f_{a,b}(n) = n^2 + an + b$. The \emph{prime chain length} $L(a,b)$ is the largest $m \geq 0$ such that $f_{a,b}(n)$ is prime for all $0 \leq n < m$.
\end{definition}

\begin{lemma}[Necessity of $b$ Prime]\label{lem:bprime}
If $L(a,b) \geq 1$, then $b$ is a positive prime.
\end{lemma}

\begin{proof}
$f_{a,b}(0) = b$ must be prime, so $b \geq 2$.
\end{proof}

\begin{lemma}[Parity Constraint]\label{lem:parity}
If $b \geq 5$ is an odd prime and $L(a,b) \geq 2$, then $a$ is odd.
\end{lemma}

\begin{proof}
$f_{a,b}(1) = 1 + a + b$ has parity equal to that of~$a$ (since $b$ is odd). If $a$ is even, then $f_{a,b}(1)$ is even, forcing $f_{a,b}(1) = 2$, i.e., $a = 1 - b$. Then $f_{a,b}(2) = 6 - b < 2$ for $b \geq 5$, contradicting $L(a,b) \geq 2$.
\end{proof}

\begin{theorem}[Connection to Heegner Numbers]\label{thm:heegner}
The discriminant of $n^2 + n + 41$ is $\Delta = -163$. By the Stark--Heegner theorem, $h(-163) = 1$, and by Rabinowitz's theorem, $n^2 + n + 41$ is prime for $n = 0, \ldots, 39$.
\end{theorem}

\begin{proof}
The class number $h(-163) = 1$ implies $\mathbb{Z}\bigl[\frac{1+\sqrt{-163}}{2}\bigr]$ is a PID. By Rabinowitz (1913), $n^2 + n + p$ produces primes for $n = 0, \ldots, p-2$ if and only if $h(1-4p) = 1$. For $p = 41$: $1 - 4(41) = -163$, confirming the result.
\end{proof}

\begin{remark}
The quadratic $n^2 - 79n + 1601 = (n-40)^2 + (n-40) + 41$ is a reindexing of Euler's polynomial.
\end{remark}

\begin{theorem}[Optimal Coefficients]\label{thm:optimal}
The pair $(a, b) = (-61, 971)$ uniquely maximizes $L(a,b) = 71$ among $|a| < 1000$, $|b| \leq 1000$. Its discriminant is $61^2 - 4 \cdot 971 = -163$.
\end{theorem}

\begin{proof}
Exhaustive search over the 168 primes $b \leq 1000$ and admissible values of~$a$ (applying Lemmas~\ref{lem:bprime}--\ref{lem:parity}) yields a maximum chain length of 71, achieved uniquely at $(-61, 971)$.
\end{proof}

\section{Algorithm}

\begin{algorithm}[H]
\caption{Quadratic Primes}
\begin{algorithmic}[1]
\REQUIRE Bounds $A_{\max}$, $B_{\max}$
\ENSURE Product $a \times b$ for optimal $(a, b)$
\STATE Compute prime sieve up to $A_{\max}^2 + A_{\max} \cdot B_{\max} + B_{\max}$
\STATE $\mathcal{B} \gets \{p \leq B_{\max} : p \text{ is prime}\}$
\STATE $\text{best} \gets 0$, $\text{result} \gets 0$
\FOR{each $b \in \mathcal{B}$}
    \FOR{$a = -(A_{\max}-1)$ \TO $A_{\max}-1$ (odd $a$ only when $b \geq 5$ is odd)}
        \STATE $n \gets 0$
        \WHILE{$n^2 + an + b \geq 2$ \AND $n^2 + an + b$ is prime}
            \STATE $n \gets n + 1$
        \ENDWHILE
        \IF{$n > \text{best}$}
            \STATE $\text{best} \gets n$, $\text{result} \gets a \times b$
        \ENDIF
    \ENDFOR
\ENDFOR
\RETURN $\text{result}$
\end{algorithmic}
\end{algorithm}

\section{Complexity Analysis}

\begin{proposition}
The algorithm runs in $O(|\mathcal{B}| \cdot |A| \cdot C)$ time and $O(S)$ space, where $|\mathcal{B}| = 168$, $|A| \leq 1000$ (after parity pruning), $C \leq 71$, and $S \approx 10^5$ is the sieve size.
\end{proposition}

\begin{proof}
The triple loop performs at most $168 \times 1000 \times 71 \approx 1.2 \times 10^7$ sieve lookups, each costing $O(1)$.
\end{proof}

\section{Answer}

\[
\boxed{-59231}
\]

\end{document}
